\chapter{Stand des Wissens}\label{cha:relatedWork}

Predictive Maintenance nutzt Zustandsdaten, um Wartungen besser zu planen. Dafür werden Verfahren benötigt, die den Zustand einer Maschine bewerten und ihre weitere Entwicklung vorhersagen. Auch der wirtschaftliche Nutzen der Wartungsentscheidung muss berücksichtigt werden \parencites{voglReviewDiagnosticPrognostic2016}{zontaPredictiveMaintenanceIndustry2020}{verbertTimelyConditionbasedMaintenance2017}. Eine wichtige Prognosegröße ist die \ac{RUL}. Sie kann mit physikalischen, datenbasierten oder hybriden Verfahren geschätzt werden. Neuere Übersichten zeigen, dass dafür zunehmend tiefe neuronale Netze eingesetzt werden. Probleme bestehen weiterhin bei der Datenqualität, der Vergleichbarkeit der Versuche und der Übertragung auf andere Anlagen \parencites{huRemainingUsefulLife2019}{berghoutSystematicGuidePredicting2022}{wuRemainingUsefulLife2024}{aminRemainingUsefulLife2026}.

\ac{CMAPSS} wird häufig als Benchmark für datenbasierte \ac{RUL}-Verfahren verwendet. Der Datensatz enthält simulierte Verläufe von Turbofan-Triebwerken mit unterschiedlichen Betriebs- und Fehlerbedingungen \parencite{saxenaDamagePropagationModeling}. Ji et al. reduzieren die Zahl der Merkmale mit einer Hauptkomponentenanalyse und verarbeiten die Zeitreihen anschließend mit einem bidirektionalen \ac{LSTM} \parencite{jiRemainingUsefulLife2020}. Li et al. verwenden stattdessen ein tiefes \ac{CNN}, das zeitliche Fenster der Sensordaten verarbeitet \parencite{liRemainingUsefulLife2018}. Listou Ellefsen et al. kombinieren unüberwachtes Vortraining mit einer überwachten \ac{RUL}-Schätzung. Dadurch kann das Modell auch mit weniger Verläufen trainiert werden, für die \ac{RUL}-Werte vorliegen \parencite{listouellefsenRemainingUsefulLife2019}. Wang et al. verbinden einen Convolutional Autoencoder mit einem \ac{TCN} und prüfen das Modell ebenfalls auf \ac{CMAPSS} \parencite[6--10]{wangRemainingUsefulLife2020}. Diese Arbeiten zeigen, dass sowohl rekurrente Modelle als auch temporale Faltungen für die \ac{RUL}-Prognose geeignet sind. Bei der Auswertung werden jedoch meist vollständige Sensorkanäle verwendet.

\acp{TSFM} werden auf großen Sammlungen unterschiedlicher Zeitreihen vortrainiert und können anschließend für neue Aufgaben verwendet werden. Chronos wandelt Messwerte in diskrete Tokens um und nutzt anschließend Sprachmodellarchitekturen für probabilistische Prognosen \parencite{ansariChronosLearningLanguage}. Das Modell \ac{TimesFM} teilt die Eingabe in Abschnitte und erzeugt mit einem Decoder Prognosen für zuvor unbekannte Datensätze \parencite{dasDecoderonlyFoundationModel2024}. MOMENT unterstützt Prognose, Imputation, Klassifikation und Anomalieerkennung. Eine \ac{RUL}-Regression wird in der ursprünglichen Untersuchung jedoch nicht betrachtet \parencite[3--5]{goswamiMOMENTFamilyOpen2024}. Dintén und Zorrilla verwenden MOMENT als unveränderten Merkmalsextraktor für \ac{CMAPSS}. Bei kleinen Trainingsmengen zeigen sich Vorteile für einzelne nachgeschaltete Modelle, jedoch nicht für alle untersuchten Verfahren. Fehlerhafte Sensordaten werden nicht untersucht \parencite[246, 252--255]{dintenUsingTimeSeries2025}. El-Ghoussani et al. verbinden Chronos-2 mit einem trainierbaren Regressionskopf und vergleichen den Ansatz mit \ac{LSTM}, \ac{GRU} und \ac{TCN}. Längere Datenlücken werden vor der Auswertung entfernt. Wie robust das Modell gegenüber solchen Fehlern ist, bleibt daher offen \parencite[2--4]{el-ghoussaniTimeSeriesFoundationModel2026}.

Fehlende Werte in Zeitreihen werden auch unabhängig von der \ac{RUL}-Prognose untersucht. \ac{GRUD} berücksichtigt dafür Masken und die Zeit seit der letzten vorhandenen Beobachtung \parencite{cheRecurrentNeuralNetworks2018}. \ac{BRITS} verarbeitet eine Zeitreihe vorwärts und rückwärts, um fehlende Werte aus früheren und späteren Beobachtungen zu schätzen \parencite{caoBRITSBidirectionalRecurrent}. Beide Verfahren wurden nicht für die \ac{RUL}-Prognose auf \ac{CMAPSS} entwickelt. Zhang und Liu verbinden dagegen die Rekonstruktion fehlender Werte direkt mit der \ac{RUL}-Schätzung in einem \ac{CNN}-\ac{LSTM}-Modell. Die Trainingsdaten werden bereits künstlich beschädigt, sodass das Modell die untersuchten Fehlermuster beim Training kennenlernt \parencite[4--6, 9--14]{zhangLSTMBasedMultiTaskMethod2023}. Tang untersucht Datenverluste, Sensorausfälle, Messabweichungen und Drift. Ein generatives Modell rekonstruiert zuerst die Sensordaten, bevor ein weiteres Modell die \ac{RUL} schätzt. Ein Vergleich mit einem \ac{TSFM} fehlt \parencite[3--6, 9--10, 19]{tangDeepLearningFramework2026}.

Neben der Prognosequalität ist der Rechenaufwand der Modelle relevant. Henderson et al. empfehlen, Rechenzeit, verwendete Hardware und Energiebedarf nachvollziehbar zu dokumentieren \parencite{hendersonSystematicReportingEnergy}. Auch in den \ac{RUL}-Studien werden solche Fragen teilweise aufgegriffen: Dintén und Zorrilla berichten Trainings- und Inferenzzeiten, während Tang den Aufwand seines Verfahrens zur \ac{RUL}-Prognose unter Sensorfehlern untersucht \parencites[255]{dintenUsingTimeSeries2025}[19]{tangDeepLearningFramework2026}. Wen und Chen vergleichen \ac{TSFM} mit kleineren Modellen. In ihrem \ac{CMAPSS}-Versuch erzielt ein \ac{TCN}-Autoencoder bessere Klassifikationsergebnisse als MOMENT und benötigt weniger Rechenzeit und belegt weniger Speicher auf der \ac{GPU}. Da Risikoklassen statt \ac{RUL}-Werten vorhergesagt werden, lässt sich das Ergebnis nicht direkt auf die \ac{RUL}-Regression übertragen \parencite[2--5]{wenTimeSeriesFoundationModels2026}.

Die bisherigen Arbeiten zeigen, dass sowohl die \ac{RUL}-Prognose mit \ac{TSFM} als auch der Umgang mit fehlerhaften Sensordaten bereits untersucht werden. Einzelne Studien berücksichtigen dabei auch den Rechenaufwand. In der bisherigen Literaturrecherche wurde jedoch kein Vergleich gefunden, der MOMENT, \ac{LSTM} und \ac{TCN} für die \ac{RUL}-Prognose mit denselben gezielt veränderten Testdaten untersucht und zugleich den Aufwand ihres Cloud-Betriebs erfasst. Offen bleibt damit, wie sich die Prognosequalität dieser Modellansätze bei zunehmenden Datenfehlern verändert und welcher Ressourcenbedarf damit verbunden ist. Die geplante Arbeit untersucht diesen Zusammenhang unter gemeinsamen Versuchsbedingungen. Ihr Beitrag liegt im Vergleich der Modelle, nicht in der Entwicklung eines neuen Prognosemodells.

\newpage
